\cleardoublepage

\chapter*{Abstract\markboth{Abstract}{Abstract}}

Home healthcare has become one of the main challenges associated with population ageing and the increasing prevalence of chronic diseases. In this context, the correct administration of medication and the periodic monitoring of physiological parameters play a key role in improving patients' quality of life. Recent advances in embedded systems and mobile technologies have enabled the development of new tools aimed at facilitating these tasks in domestic environments.

Despite these advances, many existing commercial solutions are either expensive or provide only limited functionality. In addition, medication management and biomedical monitoring are often implemented through separate devices, making it difficult to obtain a unified view of the user's health information. Consequently, there is a growing interest in developing accessible solutions capable of integrating these functionalities into a single platform.

To address this problem, this project presents \textit{MyMeds}, an intelligent home healthcare assistance system consisting of a physical device based on an ESP32 microcontroller with a touchscreen interface and an Android mobile application. The system allows users to manage medications and scheduled doses, synchronise information via WiFi, store medication records and monitor heart rate using a biomedical sensor. It also incorporates persistent storage mechanisms and PIN-based authentication to ensure both data preservation and secure access.

To validate the proposed solution, several experiments were carried out focusing on biomedical monitoring, data synchronisation, data persistence and communication security. The obtained results demonstrate the correct operation of the implemented functionalities and confirm the technical feasibility of the developed system as a support tool for medication management in home environments.
